% Transcription — fonds Alexandre Grothendieck, shelfmark 154, batch 5 (pages 81-100).
% Established from the facsimile archives/batches/154/batch-05.pdf, cut from a
% copy of 154.pdf fetched through the project's own relay
% (https://grothendieck-relay.onrender.com/source/154.pdf); PDF page k =
% archivists' page 80+k, no cover sheet. Per the skill transcribe-grothendieck.
% The folder is 175 pages in nine batches, transcribed in parallel; this is
% the fifth, and no neighbouring file was relied on. Notation follows the
% transcription of folder 80 (Systèmes de pseudo-droites).
%
% Pass: Opus 5.5 (claude-opus-5-5), 2026-09-26 — first pass, unchecked against the pages by a human.
%
% Dating and title. \dating is the folder's own date, copied verbatim from
% src/content/catalogue.ts; the folder belongs to the inventory group
% « Espaces stratifiés » (150-151), but since it has a date of its own the
% group's range is not added. Title from the catalogue, trailing full stop
% dropped; « [Système de pseudo-droites] » is bracketed there because the
% archivists supplied it.
%
% Paper. Nearly every leaf is fan-fold computer listing paper. Pages 81, 83,
% 90 and 95 were scanned upside down (printed banners of 16 August 1982 on
% 90, 95, 97, 19 August 1982 on 99): turned and read.
%
% Pages skipped: 98, a blank cream separator sheet.
% Pages summarised: none. Pages given only as a figure described in a French
% \note, with their few labels and formulas set: 82, 84, 88, 95, 97, 99, 100;
% largely figures with some text: 83, 90, 93, 94, 96.
% Pages transcribed from his hand: 81, 85-87, 89, 91-92, 93-94 (text), 96.
%
% His pagination and dates. No page carries a number of his. Two dates in his
% hand: « 5.1.84 » (underlined) at the head of p. 85 and « 5.1.84 » at the
% head of p. 87. Numbered runs: the list a) a') b) c) on p. 90; the cases
% 1) and 2°) on p. 92; the alineas marked (PS) on pp. 91 and 92; a
% small run of arrows numbered 1-5 in a drawing on p. 95. All recorded as
% written, no \note needed beyond those given.
%
% What the batch is about.
%   81     Characters chi_0 (orientation of D), chi_1 (the two inscribed
%          n-gons), chi_2 (parity) of the dihedral monodromy; tables of their
%          values on half-loops around an edge and half-turns around a point of
%          order alpha; chi_0 chi_2 is the orientability character of the
%          surface of relative positions; for n even, pi -> D_2n surjective
%          and image containing D_2n x {±1}. Cancelled by three diagonals.
%   83     A count n - (nu' - 1), 2n - [(nu - 1) + (nu' - 1)] with sketches.
%   85-86  (5.1.84) If X is the surface of non-supersingular relative
%          positions, allowing crossings of supersingular positions,
%          pi_1(X, D) -> D_2n (= Aut Pol(D)) is surjective; without crossings
%          (X~) the image contains the rotations and the reflections in sides,
%          of index <= 2; by comparing loops that cross a supersingular
%          position with half-turns, pi_1(X~) -> D_2n is surjective in all
%          cases, i.e. the principal D_2n-cover of X~ is connected.
%   87-89  (5.1.84) Generalisation: a compact surface X with a cellular graph
%          X_1 (even vertices, e.g. a pseudo-line system); circles D in
%          standard position, isotopy (same relative position), sweeping
%          (balayage) over an even vertex, specialisation/generalisation,
%          transition arcs on the double cover D~, parallel transport of the
%          half-contour structure; the case of one-sided positions, where the
%          vertices of the new surface are of order 4 and correspond to pairs
%          (D, a) with a an orientation arc.
%   90     List: a) marked banks => a') marked corners, b) critical arcs on
%          edges, c) weights of faces.
%   91-92  Building (X', K' u L') cellular; exception to « vertices of order
%          4 »; (PS): every half-circle arc of D~ contains an ordinary
%          transition arc except in two cases 1) (v in X_0) and 2°) (v not in
%          X_0), which is why the vertices of X are of two kinds only (order 4,
%          and those of type (I) (D_i, beta, omega)); otherwise no control on
%          whether parallel transport on Pol(D) comes from a local system.
%   93     Figure of the supersingular gluing; the other version flattens the
%          supersingular band onto its axis; tau_X = a^nu tau_X (a the
%          antipody, nu the number of crossings).
%   94     A map with vertices and faces of even order gives a circular
%          sequence of five commuting 2:2 involutions sigma_0 sigma_f sigma_1
%          sigma_s sigma_2, hence a pentagonal map with vertices of order 2
%          or 4; « Je n'arrive pas bien à visualiser cette dernière carte ».
%   96     Hypotheses: some pseudo-line misses s, t (not III_{p,q}), resp.
%          misses u, u' (not I_n); exceptional specialisation arc gamma of
%          length 0 when only one does.
%   82, 84, 88, 95, 97, 99, 100  Coloured drawings: discs with pseudo-line
%          arrangements and swept zig-zag bands, bands crossed by pencils,
%          trees, hemispheres.
%
% Notation fixed once for the batch: \mathcal{X} for his cursive barred X
% (the surface of relative positions, as in folder 80); X, X_1, X_0 = S(X_1)
% for the surface, its graph and vertex set; \mathbb{D}_{2n} for his
% open-faced D (dihedral group); \widetilde{D} for the double cover of a
% position D; \operatorname{Pol}(D); « p.r. » for position relative,
% « ps. dr. » for pseudo-droite(s); \underline{a} for the antipody.
%
% Readings to check first: the connective prose of pp. 85-86 and 89 (fast,
% many \ill{}); p. 91 (cancelled (PS) paragraph and the end of the page);
% p. 92 cases 1) and 2°) and the diagonal NB in the middle column; the
% letters xi / eta in the tables of p. 81; the last two lines of p. 93.

\documentclass[11pt,a4paper]{article}
% Le préambule commun à toutes les transcriptions du fonds Grothendieck.
%
% Il tient en une idée : **l'incertitude se compose, elle ne se résout pas.**
% Une transcription qui lisse un mot douteux a détruit la seule chose pour
% laquelle on la faisait — un lecteur doit pouvoir savoir, à chaque endroit,
% ce qui était sur le feuillet et ce qui a été ajouté. Les six macros qui
% suivent sont donc l'appareil critique, et non de la décoration.
%
% Les mêmes commandes sont reconnues par `scripts/render.mjs`, qui produit la
% vue de lecture du site. Une macro ajoutée ici doit l'être aussi là-bas,
% faute de quoi le rendu échouera bruyamment — ce qui est voulu : un rendu
% silencieusement faux est pire qu'un rendu absent.

\ProvidesPackage{grothendieck}[2026/08/08 Transcriptions du fonds Grothendieck]

\usepackage{amsmath,amssymb,amsthm}
\usepackage{mathrsfs}
\usepackage{stmaryrd}
% Les diagrammes commutatifs sont partout dans ce fonds : c'est la langue
% même de la géométrie algébrique de Grothendieck, et une transcription qui
% les rendrait en prose ne transcrirait rien.
\usepackage{tikz-cd}
% Français par défaut — c'est la langue des feuillets — mais l'anglais est
% chargé aussi : la traduction et le résumé sont des documents anglais, et la
% typographie française (espaces avant les deux-points) y serait une faute.
% Ces documents basculent par \selectlanguage{english} en tête de corps.
\usepackage[english,french]{babel}
\usepackage{fontspec}
\usepackage{geometry}
\geometry{a4paper,margin=25mm,marginparwidth=28mm}
\usepackage{marginnote}
\usepackage{xcolor}
% ulem, not soul. `soul`'s \st and \ul are built on a character-by-character
% scan that XeLaTeX's font machinery breaks: the document fails with an
% undefined control sequence rather than degrading. ulem does the same two jobs
% and is Unicode-engine safe. `normalem` stops it hijacking \emph, which the
% transcriptions use for Grothendieck's own emphasis.
\usepackage[normalem]{ulem}
\usepackage{hyperref}
\hypersetup{colorlinks=true,linkcolor=black,urlcolor=black,pdfborder={0 0 0}}

% --- Métadonnées du lot -----------------------------------------------------
% Reprises telles quelles par le script de rendu, qui les affiche en tête de
% la vue de lecture. Elles fixent ce que le fichier prétend couvrir : sans
% elles, une transcription flotte sans point d'attache dans un fonds de seize
% mille feuillets.

\newcommand{\folder}[1]{\gdef\grothFolder{#1}}
\newcommand{\batch}[1]{\gdef\grothBatch{#1}}
\newcommand{\leaves}[2]{\gdef\grothFirst{#1}\gdef\grothLast{#2}}
\newcommand{\foldertitle}[1]{\gdef\grothFoldertitle{#1}}
\gdef\grothFolder{?}\gdef\grothBatch{?}\gdef\grothFirst{?}\gdef\grothLast{?}\gdef\grothFoldertitle{}

% --- Le résumé --------------------------------------------------------------
%
% Il ouvre la lecture modernisée, sous le titre « Résumé », et il est composé
% en petit. Ce n'est pas un ornement : c'est le seul passage du document qui
% ne soit pas des mathématiques, et un lecteur doit pouvoir le reconnaître
% avant de l'avoir lu — puis le sauter s'il connaît déjà le sujet, ou s'y
% arrêter s'il ne le connaît pas. Le filet qui le suit marque la reprise du
% corps.

\newenvironment{resume}
  {\small\setlength{\parskip}{0.45em}}
  {\par\medskip\hrule\medskip}

% --- Mots-clés ---------------------------------------------------------------
%
% En anglais, en clôture du résumé de la lecture modernisée : le vocabulaire
% moderne sous lequel chercher ce que les feuillets construisent. Le manifeste
% du site (`npm run manifest`) les extrait pour étiqueter la cote entière —
% cette ligne est la source unique des tags, il n'y a pas de fichier à côté.
\newcommand{\keywords}[1]{\par\smallskip\noindent
  {\footnotesize\textsc{Keywords} --- \textit{#1}}\par}

% --- L'appareil critique ----------------------------------------------------

% Le numéro de feuillet, en marge. C'est l'ancre qui permet de régler un
% désaccord sur une formule en regardant *un* feuillet plutôt qu'une chemise
% de six cents. Le site s'en sert aussi pour tourner les pages du fac-similé
% quand on fait défiler la transcription.
\newcommand{\leaf}[1]{%
  \marginnote{\footnotesize\color{gray!70}\textsf{#1}}%
  \label{leaf:#1}\ignorespaces}

% Illisible : on ne devine pas. Un mot inventé qui se lit comme les autres est
% le pire résultat possible.
\newcommand{\ill}{\textcolor{red!60!black}{[\,\ldots\,]}}

% Lecture douteuse : le mot est donné, et signalé comme incertain.
\newcommand{\uncertain}[1]{\uline{#1}}

% Ajout de l'éditeur — une lettre manquante, un mot sous-entendu.
\newcommand{\add}[1]{\textcolor{blue!50!black}{[#1]}}

% Biffé par Grothendieck lui-même. Ce qu'il a rayé fait partie du document :
% une rature dit souvent où la pensée a changé de direction.
\newcommand{\struck}[1]{\sout{#1}}

% Note du transcripteur, sur l'état matériel du feuillet.
\newcommand{\note}[1]{\par\smallskip{\small\color{gray!60!black}\textit{[#1]}}\par\smallskip}

% Note marginale de Grothendieck, distincte du corps du texte.
\newcommand{\marginal}[1]{\par\smallskip\noindent\hspace{1em}%
  {\small\color{gray!60!black}#1}\par\smallskip}

% --- Titre ------------------------------------------------------------------

\AtBeginDocument{%
  \begin{center}
    {\large\bfseries Fonds Alexandre Grothendieck --- cote n\textsuperscript{o}~\grothFolder}\\[2pt]
    {\itshape\grothFoldertitle}\\[4pt]
    {\small feuillets \grothFirst--\grothLast\ (lot \grothBatch)}\\[2pt]
    {\footnotesize\color{gray!60!black}Transcription établie d'après le fac-similé numérisé
     par l'Université de Montpellier.\\
     Les lectures incertaines sont soulignées, les illisibles notées [\,\ldots\,],
     les ajouts entre crochets.}
  \end{center}
  \vspace{1em}}

\endinput


\folder{154}
\batch{5}
\pages{81}{100}
\dating{1983-1984}
\watermark{Édition de démonstration}
\foldertitle{[Système de pseudo-droites] : notes manuscrites (1983-1984, s.d.)}

\begin{document}

\page{81}
\note{calculs disposés en tableaux sur une feuille de listing d'ordinateur ; la page est traversée de trois longs traits obliques au crayon, qui semblent l'annuler}
\[
\begin{array}{ccc} \chi_0 & \chi_1 & \chi_2 \\ 1 & \bar{\varepsilon} & 0 \\ 1 & \xi & 0 \\ 1 & \eta & 1 \end{array}
\qquad
\begin{array}{ccc} 1 & 1 & 0 \\ 1 & \xi & 0 \\ 1 & \eta & 1 \end{array}
\qquad
\begin{array}{ccc} 1 & 1 & 0 \\ 1 & 0 & 1 \\ 1 & \eta & 1 \end{array}
\]
\note{trois tableaux de trois lignes ; les lettres lues $\xi$ et $\eta$ (deux signes bouclés) sont d'une lecture incertaine ; dans le troisième, les entrées de la deuxième et de la troisième ligne sont surchargées et d'une lecture incertaine. Sous le premier : « $\xi - 1$ » ; sous le troisième : « $1 - 2 = 1$ »}

$\chi_0$ orientation de $D$ ; $\chi_1$ ens.\ des \add{deux} $n$-gones inscrits ; $\chi_2$ \uncertain{ensemble} parité.
\note{les deux premières lignes sont réunies par une accolade vers « \uncertain{caractère} \uncertain{fondamental} sur $\mathbb{D}_2$ »}

\emph{NB} Le \uncertain{caractère} d'orientabilité pour $\mathcal{X}$ est $\chi_0 \chi_2$, le caractère « $n$-gones inscrits » est $\chi_0 \chi_1$.

$\varepsilon = {}$parité de $n \in \lbrace 0, 1 \rbrace$, $\bar{\varepsilon} = {}$opposé de $\varepsilon = \varepsilon - 1 \bmod 2$.

$l$ lacet en $F$ ; $u(l)$ : \uncertain{automorphisme} de $\operatorname{Pol}(F)$ (polygone \uncertain{standard} d'ordre $2n$) ;
$\chi_0(l) \overset{\mathrm{def}}{=} \operatorname{sg} u(l) = \lbrace$
\note{la formule s'arrête sur l'accolade ouvrante ; « $\operatorname{sg}$ » est \uncertain{lu} sous une surcharge, et deux petits sigles biffés précèdent $u(l)$ et $\chi_0(l)$}

Valeurs de $(\chi_0, \chi_1, \chi_2)$ :
\begin{itemize}
\item $\frac12$ lacet autour d'une arête (ou $2n+2$ pas) : $\quad 1 \quad \bar{\varepsilon} \quad 0$ ;
\item $\frac12$ tour autour d'un point d'ordre $\alpha$ : $\quad 1 \quad ? \quad \alpha$.
\end{itemize}
\note{dans la colonne de gauche, plusieurs mots sont biffés et récrits : « $\frac12$ lacet » remplace un mot biffé, « d'une arête » un autre ; sous « $\frac12$ tour autour » plusieurs lignes biffées (« \ill{} \uncertain{sommet} \ill{} d'un point \ill{} ») aboutissent à « d'ordre $\alpha$ ». À la première ligne, la seconde valeur, « $\bar{\varepsilon}$ », est écrite par-dessus un « $1$ »}

\struck{Si $\alpha$ pair alors $\chi_0$ et}
Prenant $\alpha$ \struck{pair} impair, on voit que $\chi_0$, $\chi_2$ \emph{indépendants} i.e.\ $\chi_0$, $\chi_{\mathcal{X}}$ \uncertain{indépendants} [et $\chi_1$, $\chi_2$ aussi si $n$ pair, i.e.\ $\varepsilon = 1$]
\note{« $\chi_0$, $\chi_2$ indépendants » est souligné ; $\chi_{\mathcal{X}}$ est notre lecture de l'indice, un $\mathcal{X}$ cursif comme dans le NB ci-dessus}

$\chi_{0\ldots}$, \emph{$n$ pair} : $\pi \to \mathbb{D}_{2n}$ surjectif, \uncertain{et} $\pi \to \mathbb{D}_{2n} \times \lbrace \pm 1 \rbrace$ a une image qui contient $\mathbb{D}_{2n} \times \lbrace \pm 1 \rbrace$
\note{« $\chi_{0\ldots}$ » : un $\chi$ à indice illisible, en tête de la ligne}

\begin{itemize}
\item \emph{$n$ pair} : $\quad 1 \quad 1 \quad 0$ \quad et \quad $1 \quad ? \quad \alpha$ ;
\item \emph{$n$ impair} : $\quad 1 \quad 0 \quad 0$ \quad et \quad $1 \quad ? \quad \alpha$.
\end{itemize}
En face : $\chi_2$ n'est \uncertain{prop.}\ ni à $\chi_0$, ni à $\chi_1$ --- il est possible que $\chi_0$, $\chi_1$ le soient, \uncertain{d'où} \uncertain{alors} $? = 1$ ; et, pour $n$ impair, $\chi_0$.
\note{« il est possible que » est écrit sur la ligne au-dessus et une ligne en crochet enferme le reste ; l'ordre de lecture de ces lignes de droite est incertain}

\page{82}
\note{page de dessins en couleur, sans texte : quatre disques (bords vert clair), chacun traversé d'un arrangement de pseudo-droites tracées en brun, orange, gris et bleu, certaines en pointillé ; dans les deux disques du haut, des régions triangulaires près du bord et au centre sont coloriées en orange, et, dans celui de droite, des bandes sont hachurées en bleu ; dans les deux disques du bas, une bande en zigzag (gris foncé, hachurée) traverse l'arrangement de haut en bas, et deux points opposés du bord sont marqués. À gauche, un petit croquis d'un chemin en bleu et gris coupé de segments, et quelques traits de couleur (essais de feutres)}

\page{83}
\note{en haut à gauche, un disque traversé de six droites en noir et d'une courbe orange, qui les coupe presque toutes en deux paquets près de deux points}
\[
n - \nu - (\nu - 1) - (\nu' - 1) + \nu \;=\; n - (\nu' - 1)
\]
\[
2n - \underbrace{\bigl[ (\nu - 1) + (\nu' - 1) \bigr]}
\]
\note{dans la première ligne, le premier $\nu$, le $\nu$ de $(\nu - 1)$ et le $+\nu$ final sont barrés d'un trait oblique, marques de simplification plutôt que ratures ; l'accolade sous le crochet de la seconde ligne n'a pas de légende. Au-dessous, deux croquis : un rectangle dont les deux grands côtés portent de petits traits, étiqueté $\nu'$ à gauche et $\nu$ à droite ; un segment joignant deux points, par chacun desquels passe un faisceau de droites, étiquetés $\nu' - 1$ et $\nu - 1$. Le reste de la feuille est vide}

\page{84}
\note{page de dessin, sans texte : un disque (bord vert clair) traversé d'un arrangement de courbes orange et brunes et de fines lignes grises parallèles (un balayage) ; une bande en zigzag, bordée de brun et coloriée en orange, le traverse de haut en bas, avec une ligne médiane en tirets vert foncé ; deux points opposés du bord sont marqués. C'est une version mise au propre des deux disques du bas de la page 82}

\page{85}
\note{daté en tête, souligné, de sa main : « 5.1.84 »}
J'ai \uncertain{vérifié} que si $\mathcal{X}$ désigne la surface des pos.\ rel.\ non supersingulières, sans bord en permettant le franchissement des positions supersingulières par une arête ou par un sommet (compatibilité aux \emph{poids} des \ill{}, donc que les autres opérations qui \uncertain{s'interprètent} \ill{} \ill{} par la carte, sauf les spéc.\ et gén.\ ordinaires, \uncertain{i.e.} de poids $\pm 1$) --- alors
\[
\pi_1(\mathcal{X}, D) \longrightarrow \mathbb{D}_{2n}
\]
(ou plutôt $\operatorname{Aut}\operatorname{Pol}(D_{\cdots})$)
était surjectif --- son image contient, pour une arête \ill{} un \ill{}, \uncertain{donnée} en $D$, la symétrie par \uncertain{rapport à} cette arête \ill{} sommet, \uncertain{par le} lacet \ill{} \ill{} un \emph{« demi-tour »} de $D$ autour de \struck{\ill{}} sommet, \ill{} cette arête. Ceci ne prouve pas que
\[
\pi_1(\widetilde{\mathcal{X}}, D) \longrightarrow \operatorname{Aut}\operatorname{Pol}(D) \simeq \mathbb{D}_{2n}
\]
soit surjectif, si $\widetilde{\mathcal{X}}$ désigne le demi-\ill{} de $\mathcal{X}$ \uncertain{des} \struck{\ill{}} lieux supersinguliers --- i.e.\ on exclut les franchissements des positions supersingulières. Le raisonnement montre seulement que l'image contient les symétries p.r.\ à \struck{\ill{}} \emph{côtés} (mais pas \struck{\ill{}} celles p.r.\ à un sommet). \struck{\ill{} \ill{} \ill{}} engendrent dans $\mathbb{D}_{2n}$ un \struck{\ill{}} groupe d'indice $2$, à savoir celui formé par les \struck{rotations} les \uncertain{éléments} du sous-groupe des rotations $\mathbb{D}_{2n}^{+}$ (d'ordre $n$, \uncertain{indice} $2$), et les réflexions en question. [NB. C'est un des trois sous-groupes d'indice deux, les deux autres étant le groupe engendré \struck{\ill{}} par les réflexions qui laissent fixe \uncertain{deux} sommets, et le sous-groupe $\mathbb{D}_{2n}^{+}$ des rotations.] L'image de $\pi_1(\widetilde{\mathcal{X}})$ dans $\mathbb{D}_{2n}$ est donc d'indice $1$ ou $2$.

Faisant un \struck{\ill{} tour} demi-tour \struck{complet} autour d'une position supersingulière, \uncertain{on le} franchissant, \uncertain{tout} \ill{} \ill{} \struck{\ill{} le \ill{} \ill{} rotation \ill{} \ill{} $\operatorname{Pol}(D)$ qui \ill{} $\frac{1}{2}$ \ill{} \ill{} $\Sigma(\ldots)$ ou \ill{}} \uncertain{d'une} \ill{} \ill{} $(n+1)$ --- si $n$ est \struck{pair} impair, cela montre que l'effet sur l'ensemble \struck{\ill{}} $\Sigma_a(D)$ des \struck{\ill{}} \uncertain{polygones} \uncertain{circonscrits} (\uncertain{choisis} sur les deux autres \ill{} : deux éléments $\Sigma_s(D)$, $\underline{\omega}(D)$) \ill{} \struck{\ill{}} \ill{} pour voir si \ill{} \ill{} en $D_0$ \struck{fixé} donné \struck{fixe}
\note{« $n+1$ » est entouré, et un trait le relie à « impair » au-dessous, qui remplace « pair » biffé. Le passage de quatre lignes qui précède est biffé de traits horizontaux ; notre lecture de la syntaxe de tout ce paragraphe est très incertaine}
\marginal{\ill{} \ill{} par un \uncertain{sommet} \ill{} \ill{} \ill{} \ill{} \ill{} \ill{} \ill{} \ill{} \ill{} \uncertain{les} \ill{} \ill{} \ill{} \ill{} \ill{} \uncertain{circonstances}\ldots\ ne suffit-il \ill{} l'\uncertain{explicitation} \ill{} \ill{} un fait (\ill{}) mais l'\ill{} \ill{} \ill{} \uncertain{possible}}
\note{note écrite en biais dans l'espace entre les deux colonnes, en partie biffée de quatre traits ; nous en lisons quelques mots : « \ill{} \ill{} \ill{} \uncertain{moyennant} \ill{} \uncertain{l'explicitation} \ill{} \ill{} \ill{} \uncertain{les} \uncertain{recollements} \ill{} \uncertain{de} \uncertain{franchissements} \ill{} \uncertain{exceptionnels} », et plus bas « \uncertain{formule} \ill{} \ill{} \ill{} \ill{} un fait » avec une petite croix}

\page{86}
induit l'identité sur $\Sigma_a(D_0)$, il revient au \ill{} de se permettre de franchir les positions supersingulières, ou non : dans \uncertain{le} cas $n$ \struck{im}pair, $\pi_1(\widetilde{\mathcal{X}}) \to \mathbb{D}_{2n}$ est aussi surjectif. Si $n$ est pair, le résultat \struck{\ill{}} d'un lacet \uncertain{sur} $\Sigma_a(D_0)$ \struck{\uncertain{autour}} dans $\mathcal{X}$ qui franchit $n$ fois une position supersingulière, \uncertain{sachant} que $n$ \uncertain{est} pair \ill{}, est\ill{}, \ill{} \ill{} \uncertain{rapport} opposé à celui qu'induit \uncertain{par} \ill{} lacet $l'$ qui se déduit de $l$ en remplaçant les franchissements par des demi-tours. Si donc on \struck{\uncertain{prend}} part d'une c.\ pos.\ rel.\ $D_0$ \uncertain{que} \ill{} \ill{} \uncertain{somment} $\Sigma$ \emph{d'ordre} \uncertain{pair} \uncertain{par} \uncertain{un} \ill{}, \uncertain{s'il} \uncertain{est} \uncertain{pair}, dans ce \uncertain{mode} \ill{} les lacets de demi-tours autour de $s_0$ \add{de $D_0$}, de la façon indiquée, on trouve un \struck{\ill{}} contour \struck{de} $\Sigma(D_0)$ qui \uncertain{n'est} \uncertain{pas} l'identité. Cela prouve donc que \emph{dans tous les cas}, $\pi_1(\widetilde{\mathcal{X}}) \to \mathbb{D}_{2n}$ est surjectif --- i.e.\ le \uncertain{revêtement} \add{principal correspondant} de $\widetilde{\mathcal{X}}$, de groupe $\mathbb{D}_{2n}$, est \emph{connexe}.
\note{« principal correspondant » est ajouté au-dessus de la ligne. Au bas de la page, à gauche, une note isolée : « $\mathbb{D}_{2n}$ »}

\page{87}
\note{daté en tête : « 5.1.84 »}
Soit $X$ une surface compacte, munie d'un graphe topologique $X_1 \subset X$ la décomposant cellulairement. [On s'intéressera plus particulièrement au cas où tous les \struck{\ill{}} sommets de $X_1$ sont d'ordre pair, p.\ ex.\ $(X, X_1)$ est un syst.\ de pseudo-droites dans un plan \struck{\ill{}} projectif réel topologique.]
\struck{Soit}. Un cercle topologique $D \subset X$ sera dit « en position \uncertain{standard} » (p.r.\ à $X_1$), s'il \struck{\ill{}} rencontre $X_1$, et si pour tout $s \in X_1 \cap D$, \struck{\ill{} coupe \ill{}}, $s$ \uncertain{est} un pt d'ordre pair de $X_1$, et $D$ \emph{coupe} toutes les branches de $X_1$ en $s$. On dit que deux tels $D$, $D'$ sont « en même position relative » \ill{} \ill{} sont \emph{isotopes}, si on peut passer de $D$ à $D'$ par une « isotopie », au cours de laquelle tous les positions intermédiaires $D_t$ sont en position st.\ mais où les $D_t \cap S(K)$ \uncertain{restent} constants.
\note{au-dessus de « si on peut passer », entre parenthèses : « (\uncertain{pour} \uncertain{des} \uncertain{paramétrisations} \uncertain{continues} $D \simeq \mathbf{U}$, $D' \simeq \mathbf{U}$) », où $\mathbf{U}$ rend son U double}
Il faut dire pour cela que, \struck{\ill{}} si l'on désigne par $\Phi(D)$ l'ens.\ des facettes de $(X, X_1)$ rencontrées par $D$, que l'on ait $\Phi(D) = \Phi(D')$. Il faut, aussi qu'il existe une famille à un paramètre $f_t$ $(0 \leqslant t \leqslant 1)$ d'automorphismes de $(X, X_1)$, avec $f_0 = \mathrm{id}$, et $f_1(D) = D'$.

On définit une notion de \emph{balayage} d'une p.r.\ au-dessus d'un sommet \emph{pair} de $X_1$, comme on \ill{} dans le cas des syst.\ de ps.\ dr., et une notion correspondante de spécialisation \uncertain{vers} un sommet. \struck{Également la notion de}
Pour $D$ donné, les sp.\ et spéc.\ \uncertain{et} de générisation à partir de $D$ se répartissent en \ill{} \ill{} des « \ill{} de \struck{balayage} transition », \uncertain{sur} \ill{} dans le rev.\ $\widetilde{D}$ \uncertain{au-dessus} de $D$, qui

\note{suit, au milieu de la page, un dessin en couleur : une droite horizontale (jaune-vert, la position $D$) coupée de faisceaux de droites rouges ; à gauche un faisceau de trois droites concourantes sur $D$, le point entouré d'un petit cercle pointillé et marqué de deux petits segments sombres de part et d'autre ; puis des couples de droites qui se croisent au-dessus ou au-dessous de $D$, les segments de $D$ entre leurs traces marqués en sombre}

forment, pour sa \uncertain{variation}, une famille d'ens.\ \struck{\ill{}} \add{\uncertain{deux à deux}} disjoints. Il y a une notion de transport parallèle pour $\widetilde{D}$, muni de sa structure de \struck{\ill{}} \add{demi-}contours \uncertain{naturelle} (ou plutôt, \uncertain{sa} \uncertain{structure} combinatoire) pour les générisations et spéc.\ ordinaires, et aussi

\page{88}
\note{la page ne porte qu'un croquis, sans aucune légende : au crayon, un ovale et plusieurs courbes qui se croisent ; en vert, deux courbes issues d'un même point, fléchées de petites marques ; rien n'y est écrit}

\page{89}
Pour les familles à $1$ paramètre de $D$ qui \uncertain{décrivent} un segment \uncertain{telles} que à l'intérieur \struck{du} \uncertain{chemin} \uncertain{d'eux} la p.r.\ reste constante \add{et} \struck{\ill{}} que les extrémités \uncertain{en} \uncertain{sont} \ill{} \uncertain{deux} des $D_t$ qui restent \struck{en} en position standard).

On va s'intéresser maintenant au cas des p.r.\ à une seule \uncertain{rive} --- i.e.\ telles que $\widetilde{D}$ soit un \ill{} top. Alors l'ens.\ des sp.\ de transition \ill{} \uncertain{partant} de $D$ est muni d'une structure polygonale \add{(\uncertain{de} \uncertain{même} \uncertain{ordre} \uncertain{que} \ill{} \ill{} \uncertain{angulaires} \uncertain{critiques})} \add{\ill{} \ill{} \ill{} \uncertain{conventions} \uncertain{pour} \uncertain{orienter}} $D$ \ill{} le graphe \ill{}, et \struck{\ill{}} l'ens.\ des p.r.\ est polygonal. \ill{} définir des sp.\ de deux \uncertain{types} de \uncertain{définition} \uncertain{du} transport par l'orientation. La construction la plus intrinsèque semble celle où le \ill{} d'une généralisation « tire » sur $D$ \uncertain{pour} généraliser, et que le transport de l'orientation, \uncertain{soit} $\mathcal{X}$ \uncertain{soit} opposé au transport parallèle de l'orientation.
\note{les deux ajouts entre crochets sont écrits au-dessous et au-dessus de la ligne « de $D$ est muni d'une structure polygonale », l'un renvoyé par un trait ; leur lecture est très incertaine}

Dans ce cas, dans le contexte \uncertain{standard}, tous les sommets sont d'ordre $4$, et ils correspondent \struck{\ill{}} $1$-$1$ aux \struck{\uncertain{sympt.}\ $D$} $(D, \struck{\ill{}} a)$, où $D$ est une p.r., et $a$ un « arc d'orientation » sur $\widetilde{D}$, \uncertain{des} \uncertain{types} \uncertain{généraux} i.e.\ joignant deux \struck{sommets} \add{\ill{}} de $D \cap S(X_1)$ \uncertain{telles} que ses \ill{} \ill{} \ill{} \uncertain{ne} \ill{} ni \ill{} \uncertain{à} $S(X_1)$, ni d'autre \ill{} \ill{} \uncertain{spécialisation}).
\note{ici et plus bas il écrit $S(X_1)$ avec un $X$ qu'on pourrait aussi lire $K$ ; à la page 87, $D_t \cap S(K)$ se lit de même}

La situation est surtout symétrique si la \uncertain{décomposition} de $X$ suivant la position relative \uncertain{considérée} est un disque (ce qui signifie aussi que $X$ est un plan projectif réel, et $D$ une ps.\ droite) --- condition qui est \uncertain{invariante} par \uncertain{opérations} élémentaires, \uncertain{car} \uncertain{elles} \uncertain{peuvent} s'interpréter \uncertain{comme} \uncertain{déduites} \uncertain{de} \uncertain{ces} décompositions par des \uncertain{recollements}, et on \uncertain{récupère} des
\note{la page s'arrête sur « des »}

\page{90}
\note{feuille de listing d'ordinateur (bandeau imprimé « GUTS », daté par la machine du 16 août 1982). En haut, un croquis en couleur : une droite jaune-vert coupée par des droites orange (une paire qui se croise au-dessus d'elle, une droite seule, une autre paire) ; aucune légende}
\begin{itemize}
\item[a)] \uncertain{rives} \uncertain{marquées} $\Rightarrow$ a') \struck{angles} \add{coins} marqués
\item[b)] arcs critiques marqués sur les arêtes
\item[c)] \emph{poids} des faces
\end{itemize}
\note{une flèche double, dans la marge gauche, relie a) et c), avec une petite flèche montante à la hauteur de a)}

\page{91}
$\mathcal{X}$ un graphe topologique $L$ \uncertain{comme} à l'accoutumée. Mais dans \ill{} que les $\operatorname{Déc}(X, D)$ ne seraient pas des disques, et \uncertain{sans} \ill{} \uncertain{supposer} $\widetilde{D}$ connexe, on peut construire une \uncertain{nouvelle} $\mathcal{X}'$, \add{\uncertain{avec} \uncertain{un} \uncertain{graphe} \uncertain{top.}\ $K'$ dans $\mathcal{X}$} \uncertain{sans} \uncertain{faces} (\uncertain{on} \uncertain{oublie} \uncertain{toutes} les décomp.) --- mais $(\mathcal{X}', K')$ n'est plus \uncertain{nécessairement} cellulaire, il faut y joindre $L'$ et prendre $(\mathcal{X}', K' \cup L')$ pour avoir quelque chose de cellulaire.
\note{deux traits verticaux dans la marge gauche marquent l'alinéa qui suit}
Tantôt en disant que les sommets \add{de $(X, K)$} sont d'ordre $4$, j'ai été un peu \uncertain{rapide} --- il y a une exception, quand on \add{part d'une} \struck{\ill{}} situation \struck{\ill{}} \ill{}, où \struck{\ill{}} tout l'arc $\widetilde{uv}$ (pointillé) sur $\widetilde{D}_x$ est un arc d'orientation (position relative \emph{stable} \uncertain{donc}). La circonstance \uncertain{spéciale} dans le cas où $(X, X_1)$ est un syst.\ de ps.\ droites, et qui permet le contrôle \uncertain{sur} les sommets de $\mathcal{X}$, est la suivante

\note{dessin en couleur au milieu du bas de page : deux paires de droites rouges qui se croisent de part et d'autre d'une droite horizontale jaune-vert ($D$), prolongée en pointillé ; les points marqués $u$, $s$, $s'$, $v$ sur $D$ et $\alpha$, $\beta$ sur les droites rouges}

(PS) Soit $D$ une p.r.\ \struck{\ill{}} \struck{[\ill{} \ill{} \ill{} \ill{} \ill{} des ps.\ dr., soit $s \in D \cap X_1$,} et considérons \struck{\ill{} \ill{}} \struck{$s'$, $s''$ \ill{}} \add{antipodiques} sur $\widetilde{D}$, et une des arcs \uncertain{déterminés} \struck{\ill{}} par $s'$, $s''$ (\ill{} \ill{} un \emph{demi-cercle} \uncertain{de} polygone top.\ \uncertain{ou} comb.\ \uncertain{associé} à $\operatorname{Pol}(D)$), alors \struck{\ill{} \ill{} \ill{} des \ill{}} \struck{\ill{}} \struck{il n'y a pas d'arc de transition (« ordinaire ») dans $s \in X_0 = S(X_1)$}
\note{le début de cet alinéa (PS) est biffé : un long trait horizontal par ligne, un trait vertical à gauche, et deux grands traits obliques au crayon qui traversent toute la colonne de droite jusqu'au bas de la page ; nous ne savons pas si ces derniers annulent la suite}
Soit $a$ un arc de transition \add{(ordinaire)} sur $\widetilde{D}$, $a'$ l'arc opposé (\uncertain{donc} $\bar{a} \cap a' = \emptyset$) et considérons \struck{les} un des \uncertain{arcs} \add{$b$}\uncertain{compl.}\ \uncertain{contenu} dans $\widetilde{D} \setminus \lbrace a \cup a' \rbrace$. Alors, ou bien $b$ contient un arc de transition \add{(ordinaire)}, ou bien $b$ \struck{\ill{}} est un « arc de transition exceptionnel », i.e.\ on est dans l'un des cas suivants : on a un cercle « composant » $D_i$ de $X_1$ qui \add{a \ill{}} \uncertain{soit} sur $D_i$, une \struck{\ill{}} \add{\uncertain{facette} $\varphi$} (\uncertain{comp.}\ \uncertain{de} \add{dit \uncertain{sommets}} \ill{}) \add{\uncertain{de}} $X$ \uncertain{en} $\varphi$ \struck{\ill{} \ill{}}, une orientation \add{de} $(D_i, \varphi, \omega)$ pour et $D$ est déduit de \uncertain{déplacement} élémentaire ---
\note{toute cette fin de page est d'une lecture très incertaine}

\page{92}
(\emph{PS}) Considérons un arc-demi-cercle $a$ sur $\widetilde{D}$ ($D$ une p.r.\ \emph{de ps.\ droites}). Alors \struck{\ill{}} \add{$\mathring{a}$} contient un arc de transition (ordinaire), sauf dans les \add{deux} cas suivants :
\marginal{i.e.\ $a$ contient un arc de transition qui n'a pas d'extrémité \ill{} \ill{}}

\struck{Les \ill{}} 1) Les extrémités de $s, s'$ de $a$ \struck{\ill{}} sont \uncertain{éléments} d'arcs de généralisation \add{$\alpha$, $\alpha'$} (non \uncertain{ortho}\ill{}) [\uncertain{i.e.}\ les \uncertain{pts} $\sigma$ de $X$ définis par \add{$s$, $s'$} sont \uncertain{des} sommets de $X_1$] et il y a un \uncertain{composant} $D_i$ de $X_1$ \uncertain{passant} \uncertain{par} $\sigma$, et une orientation \add{de} $X$ en $\sigma$, telle que $D$ soit \struck{déduit de} « voisin de $D_i$ » et « déduit de $D_i$ par \add{petite} rotation \uncertain{autour} de $\sigma$ » ; $a \setminus (\mathring{\alpha} \cap a) \setminus (\mathring{\alpha}' \cap a)$ est la rive de balayage de $D$ correspondant à $(D_i, \sigma, \omega)$.
\note{une accolade à droite de ce cas 1) porte « cas $v \in X_0$ »}

2°) \struck{\ill{} \ill{}} \struck{\uncertain{les} extrémités $s, s'$ de $a$, \ill{}} \add{les} extrémités \add{de} \struck{\ill{}} \struck{\ill{}} \ill{} de spécialisation $\alpha_3$, $\alpha_5$ (\ill{} \uncertain{antipodiques} \add{\uncertain{continus} dans $a$} de \uncertain{l'autre} !) \struck{tels qu'il existe $\sigma$, $\sigma'$} \struck{Soit} \add{Soit $\beta$} \add{\uncertain{bien} \uncertain{plus}} l'arête de $X_1$ \uncertain{passant} par $\sigma$, \struck{soit} $D_i$ le \uncertain{sommet} de $X_1$ contenant $\beta$, il existe une orientation \struck{de $X_1$} \struck{\uncertain{sur}} \uncertain{de} $\beta$ \uncertain{telle} \uncertain{que} $D$ \uncertain{soit} déduite ---
\note{une accolade à droite de ce cas 2°) porte « cas $v \notin X_0$ »}

C'est cette circonstance qui \uncertain{m'assure} qu'il n'y a que deux sortes de sommets dans $\mathcal{X}$, les sommets d'ordre $4$ dont je viens de parler, et les sommets correspondant aux \uncertain{positions} \add{$(\mathrm{I})$} $(D_i, \beta, \omega)$, où $D_i$ est un \uncertain{composant} irr.\ de $X_1$, $\beta$ une \uncertain{arête} \add{\uncertain{sommet}} dans $X$, $\beta$ une facette de $D_i$ \ill{} \ill{}, $\omega$ une orientation de $X$ en $\beta$. Sinon, on n'a pas de contrôle, a priori, sur le type des « circonstances » \uncertain{sommets} sur $\mathcal{X}$ --- et il s'ensuit qu'il n'est pas clair \struck{\ill{}} dans quelle mesure le transport parallèle sur les $\operatorname{Pol}(D)$ provient d'un syst.\ local sur $\mathcal{X}$.
\note{la colonne du milieu porte, écrite en biais, une note annotée « NB » avec un petit dessin en couleur (une droite jaune et deux croix rouges marquées $T_1$, $T_2$) ; on y lit « NB Dans ce cas \ill{} \ill{} \uncertain{nous} \uncertain{arrive} \ill{} » puis une liste « a) $s$, \ill{} \uncertain{antipodiques} b) $a \neq$ \ill{} (\uncertain{antipodaux}) c) dist.\ \ill{} d) \ill{} \uncertain{spécialisations} \ill{} », une mention « \uncertain{sommets} $(\mathrm{I})$ » dans une ellipse, et, plus bas, quatre lignes biffées de traits obliques : « \ill{} \ill{} \ill{} $X$ \ill{} \ill{} droite \ill{} \ill{} \ill{} \uncertain{identiques} \ill{} \ill{} \uncertain{p.r.}\ \ill{} ». Plus bas dans la même colonne : « NB Dans ce cas, \uncertain{contient} deux arcs de spécialisation, \struck{\ill{}} (mais à \uncertain{rencontrer} \ill{}) »}

\page{93}
\note{la moitié supérieure de la page est un grand dessin en couleur : deux droites horizontales vert clair, et entre elles une ligne médiane en tirets rouges marquée de petits traits bleus ; quatre droites verticales vert clair ; des droites orange qui, entre les verticales, se croisent sur la ligne médiane en forme de sabliers ou de losanges, et deux grands arcs orange, en haut et en bas, qui se croisent sur les verticales. Aux croisements du haut, trois petites marques étiquetées $\nu$, $\nu'$ et $\nu'$, avec des arcs en tirets vert clair. Deux légendes sur la droite horizontale du haut : « $n - \nu$ » et « $n - 1 - (\nu - 1) - (\nu' - 1)$ »}

Version du recollement supersingulier adaptée à l'étude des \uncertain{applications} $\struck{\mathcal{X}} \to X$, mais non \struck{\ill{} \ill{}} des structures \uncertain{\emph{n-gonales}} des faces et à l'étude du transport parallèle de celles-ci.

Dans l'autre version, on aplatit la bande supersingulière sur son axe central (pointillé rouge). Le transport parallèle \add{(de franchissement d'une position supersingulière)} n'est \emph{pas} celui donné par l'identification des côtés utilisée pour construire $\mathcal{X}$, mais est l'antipodique de celui-ci. Dans la \uncertain{relation} entre transport parallèle sur $X$, et transport-réflexion \add{$\tau_{\mathcal{X}}$} sur $\mathcal{X}$, pour un chemin donné sur $\mathcal{X}$, \struck{parallèle \ill{} de}
\[
\tau_X = \underline{a}^{\nu} \tau_{\mathcal{X}} ,
\]
où $\underline{a}$ est l'antipodisme, et $\nu$ le \uncertain{nombre} de franchissements d'une ligne supersingulière. \uncertain{Les} \uncertain{deux} \uncertain{systèmes} \uncertain{locaux} \uncertain{orientés} \ill{} \ill{} $\mathcal{X}$, \ill{} \ill{} \ill{} \ill{} \ill{} \ill{} \ill{} \uncertain{communs} \ill{} \ill{}
\note{« $\tau_X$ » est ajouté au-dessus de « relation » et renvoyé par un trait, « $\tau_{\mathcal{X}}$ » au-dessus de « transport-réflexion » ; l'antipodisme est souligné, comme dans le dossier 80 ($\underline{a}$). Les deux dernières lignes, serrées au bas de la feuille, ne se lisent qu'en partie}

\page{94}
\note{en haut à gauche, au crayon, un croquis : des droites concourantes en un point, trois petits triangles marqués $\sigma_0$, $\sigma_1$, $\sigma_2$ ; au-dessous, une accolade (un arc) sur la suite « $\sigma_0\ \sigma_f\ \sigma_1\ \sigma_s\ \sigma_2$ », chaque terme surmonté d'un point, les indices $f$ et $s$ surchargés et un sixième terme raturé. Au centre, deux dessins en couleur de la même carte (arêtes vert clair et bleues) : un sommet où se rencontrent des arêtes des deux couleurs, un triangle colorié en orange et les secteurs voisins marqués $\sigma_0$, $\sigma_1$, $\sigma_2$, $\sigma_f$, $\sigma_s$, hachurés en gris dans le second. À droite, une ligne brisée vert clair marquée de petits traits rouges et de deux points d'exclamation, un pentagone au crayon, et, en bas, un faisceau de droites grises aboutissant à une courbe vert clair marquée de traits rouges, étiquetée « $\sigma_s$ »}
Une carte dont les sommets \emph{et} les faces sont d'ordre pair, donne naissance à une suite circulaire de cinq involutions $2 : 2$ commutantes
\[
\sigma_0\ \sigma_f\ \sigma_1\ \sigma_s\ \sigma_2\ [\sigma_0 \ldots]
\]
donc à une carte pentagonale \uncertain{frondière}, dont tous les sommets sont d'ordre deux ou quatre. Je n'arrive pas bien à \uncertain{visualiser} cette dernière carte, en termes de la carte initiale\ldots
\note{à droite, isolé et encadré de deux traits : « $\sigma_s\ \sigma_a\ \sigma_r$ », les deux premiers réunis par une accolade, et au-dessous « $\sigma_{\mathrm{vert}}$ »}

\page{95}
\note{croquis en couleur sans phrase, sur une feuille de listing d'ordinateur (bandeau imprimé du 16 août 1982). En haut à droite, deux dessins d'une droite vert clair orientée (flèche) coupée de trois paires de droites orange qui se croisent au-dessus ou sur elle ; des segments de la droite verte sont renforcés en noir, et de fines flèches au crayon indiquent des déplacements ; sous le premier croisement, une suite de petites flèches alternées numérotées $1$, $2$, $3$, $4$, $5$. Au-dessous, la légende « $\sigma_1 \quad \sigma \quad \sigma_2\ \sigma_0$ », avec une accolade étiquetée « $2n$ » sur les trois derniers termes. En bas, un petit graphe vert clair (deux sommets reliés par une arête, chacun d'où partent d'autres arêtes), avec des segments renforcés en noir sur les arêtes}

\page{96}
\note{dessins en couleur : en haut à gauche, un disque traversé de pseudo-droites orange, avec, sur le bord, les points $u$, $u'$, $s$, $t$ et les arcs marqués $\alpha$, $\beta$, $\gamma$, chaque segment de bord entre eux marqué « $\frac12$ » par une accolade ; au centre, un petit disque traversé de pseudo-droites orange ; en haut à droite, un disque de même facture, avec $u$, $u'$, des arcs $\alpha$, $\beta$ très courts et un long arc $\gamma$ en bas, les segments de bord marqués « $\frac12$ » ; en bas, deux bandes horizontales (bords vert clair, axe orange) coupées de faisceaux de droites orange, les bords renforcés en noir par segments, en tirets ou pleins, qui se croisent au milieu en un point marqué $u$, $u'$ dans la seconde}

\emph{NB} On suppose qu'il y a au moins une ps.\ dr.\ qui ne passe pas par $s$, $t$ i.e.\ qu'on n'est pas dans le cas $\mathrm{III}_{p,q}$ ($p, q \geqslant 1$ --- dans pour \ill{} le cas $\struck{\mathrm{II}_3}$ \uncertain{non} plus). Quand il n'y a qu'une seule ps.\ droite qui ne passe pas par $s$, $t$, \uncertain{l'arc} de spéc.\ exceptionnel $\gamma$ est de long.\ $0$.
\note{« $\mathrm{III}_{p,q}$ » est surligné ; le « $\mathrm{II}_3$ » de la ligne suivante est surchargé et d'une lecture incertaine}

On suppose qu'il y a une ps.\ droite qui ne passe pas par $u$, $u'$ i.e.\ qu'on n'est pas dans le cas du type $\mathrm{I}_n$. Dans le cas où il n'y en a qu'une, l'arc de spécialisation \struck{de pas} exceptionnel $\gamma$ est de longueur $0$.

\page{97}
\note{croquis en couleur sans phrase, sur une feuille de listing d'ordinateur (bandeau imprimé du 16 août 1982) : à droite, un disque traversé de pseudo-droites orange, deux petits arcs du bord renforcés en noir aux extrémités d'une corde, et un long arc noir en bas ; au milieu en bas, une bande (bords vert clair et noir, pleins ou en tirets, hachurée par endroits) coupée de droites orange, dont les bords se croisent en un point ; à droite en bas, un grand disque où un faisceau de pseudo-droites orange se tresse vers le haut, avec deux arcs du bord et un arc du bas renforcés en noir et des segments de bord marqués « $\frac12$ » ; enfin deux petits disques du même type, dans le second les segments de bord marqués « $\frac12$ », « $0$ », « $\frac12$ »}

\page{99}
\note{croquis en couleur sans texte, sur une feuille de listing d'ordinateur (bandeau imprimé du 19 août 1982) : un rectangle vert clair aux coins arrondis, rempli d'arcs orange et coupé par une bande verticale bordée de tirets gris-bleu, avec un point d'interrogation au-dessus et un au-dessous ; un arbre vert clair ramifié portant trois petits disques dans lesquels se croisent des segments orange et bruns ; un rectangle orange partagé en deux par un trait vertical}

\page{100}
\note{page de croquis en couleur, sans texte : deux arbres ramifiés vert clair ; une droite vert clair coupée de droites orange ; un rectangle vert clair partagé en trois bandes verticales (celle du milieu bordée de gris-bleu), traversé d'arcs orange qui délimitent de petits triangles coloriés en orange ; des disques et des demi-disques, bordés de vert clair ou de vert foncé et traversés de pseudo-droites orange, certains avec un arc en tirets gris-bleu ou vert (l'équateur d'un hémisphère vu en perspective) ; une surface allongée, faite d'un disque surmonté et prolongé de deux calottes vert clair, dans laquelle se tressent des courbes orange ; un rectangle aux bords orange traversé de courbes ; un carré vert clair coupé d'une bande en tirets gris-bleu où serpentent des courbes orange}

\end{document}
